\documentclass[10pt,a4paper]{amsart}
\usepackage{amsmath,amssymb,amsthm,graphicx}
% For dealing with references we use the comment environment
\usepackage{verbatim}
\newenvironment{reference}{\comment}{\endcomment}
%\newenvironment{reference}{}{}
\newenvironment{slogan}{\comment}{\endcomment}
\newenvironment{history}{\comment}{\endcomment}
% For commutative diagrams we use Xy-pic
\usepackage[all]{xy}

% We use 2cell for 2-commutative diagrams.
\xyoption{2cell}
\UseAllTwocells
\theoremstyle{plain}
\newtheorem{theorem}[subsection]{Theorem}
\newtheorem{proposition}[subsection]{Proposition}
\newtheorem{lemma}[subsection]{Lemma}

\theoremstyle{definition}
\newtheorem{definition}[subsection]{Definition}
\newtheorem{example}[subsection]{Example}
\newtheorem{exercise}[subsection]{Exercise}
\newtheorem{situation}[subsection]{Situation}

\theoremstyle{remark}
\newtheorem{remark}[subsection]{Remark}
\newtheorem{remarks}[subsection]{Remarks}

\numberwithin{equation}{subsection}
% Macros
%
\def\lim{\mathop{\mathrm{lim}}\nolimits}
\def\colim{\mathop{\mathrm{colim}}\nolimits}
\def\Spec{\mathop{\mathrm{Spec}}}
\def\Hom{\mathop{\mathrm{Hom}}\nolimits}
\def\Ext{\mathop{\mathrm{Ext}}\nolimits}
\def\SheafHom{\mathop{\mathcal{H}\!\mathit{om}}\nolimits}
\def\SheafExt{\mathop{\mathcal{E}\!\mathit{xt}}\nolimits}
\def\Sch{\mathit{Sch}}
\def\Mor{\mathop{\mathrm{Mor}}\nolimits}
\def\Ob{\mathop{\mathrm{Ob}}\nolimits}
\def\Sh{\mathop{\mathit{Sh}}\nolimits}
\def\NL{\mathop{N\!L}\nolimits}
\def\CH{\mathop{\mathrm{CH}}\nolimits}
\def\proetale{{pro\text{-}\acute{e}tale}}
\def\etale{{\acute{e}tale}}
\def\QCoh{\mathit{QCoh}}
\def\Ker{\mathop{\mathrm{Ker}}}
\def\Im{\mathop{\mathrm{Im}}}
\def\Coker{\mathop{\mathrm{Coker}}}
\def\Coim{\mathop{\mathrm{Coim}}}

% Boxtimes
%
\DeclareMathSymbol{\boxtimes}{\mathbin}{AMSa}{"02}

%
% Macros for moduli stacks/spaces
%
\def\QCohstack{\mathcal{QC}\!\mathit{oh}}
\def\Cohstack{\mathcal{C}\!\mathit{oh}}
\def\Spacesstack{\mathcal{S}\!\mathit{paces}}
\def\Quotfunctor{\mathrm{Quot}}
\def\Hilbfunctor{\mathrm{Hilb}}
\def\Curvesstack{\mathcal{C}\!\mathit{urves}}
\def\Polarizedstack{\mathcal{P}\!\mathit{olarized}}
\def\Complexesstack{\mathcal{C}\!\mathit{omplexes}}
% \Pic is the operator that assigns to X its picard group, usage \Pic(X)
% \Picardstack_{X/B} denotes the Picard stack of X over B
% \Picardfunctor_{X/B} denotes the Picard functor of X over B
\def\Pic{\mathop{\mathrm{Pic}}\nolimits}
\def\Picardstack{\mathcal{P}\!\mathit{ic}}
\def\Picardfunctor{\mathrm{Pic}}
\def\Deformationcategory{\mathcal{D}\!\mathit{ef}}
\usepackage{cleveref}
\crefname{originalequation}{original equation}{original equations}
\Crefname{originalequation}{Original equation}{Original equations}
\makeatletter
\newlabel{equation-descent-etale}{{0BTI}{source}{}{sourcepointer-1}{}}
\newlabel{equation-descent-etale@cref}{{[originalequation][0][]0BTI}{[1][0][]source}{}{}{}}
\makeatother
\crefname{originallemma}{original lemma}{original lemmas}
\Crefname{originallemma}{Original lemma}{Original lemmas}
\makeatletter
\newlabel{lemma-faithful}{{0BTJ}{source}{}{sourcepointer-2}{}}
\newlabel{lemma-faithful@cref}{{[originallemma][0][]0BTJ}{[1][0][]source}{}{}{}}
\makeatother
\crefname{originallemma}{original lemma}{original lemmas}
\Crefname{originallemma}{Original lemma}{Original lemmas}
\makeatletter
\newlabel{schemes-lemma-compose-after-separated}{{01KV}{source}{}{sourcepointer-3}{}}
\newlabel{schemes-lemma-compose-after-separated@cref}{{[originallemma][0][]01KV}{[1][0][]source}{}{}{}}
\makeatother
\crefname{originallemma}{original lemma}{original lemmas}
\Crefname{originallemma}{Original lemma}{Original lemmas}
\makeatletter
\newlabel{schemes-lemma-semi-diagonal}{{01KS}{source}{}{sourcepointer-4}{}}
\newlabel{schemes-lemma-semi-diagonal@cref}{{[originallemma][0][]01KS}{[1][0][]source}{}{}{}}
\makeatother
\crefname{originalproposition}{original proposition}{original propositions}
\Crefname{originalproposition}{Original proposition}{Original propositions}
\makeatletter
\newlabel{proposition-properties-sections}{{024T}{source}{}{sourcepointer-5}{}}
\newlabel{proposition-properties-sections@cref}{{[originalproposition][0][]024T}{[1][0][]source}{}{}{}}
\makeatother
\crefname{originaltheorem}{original theorem}{original theorems}
\Crefname{originaltheorem}{Original theorem}{Original theorems}
\makeatletter
\newlabel{theorem-etale-radicial-open}{{025G}{source}{}{sourcepointer-6}{}}
\newlabel{theorem-etale-radicial-open@cref}{{[originaltheorem][0][]025G}{[1][0][]source}{}{}{}}
\makeatother
\title{Lemma fully faithful in etale.tex}
\author{The Stacks Project Authors}
\date{stacks-project:a04446e57ec1fbc252a871afcec7752fb2807b14; GNU Free Documentation License 1.2 (as published in the Stacks Project COPYING)}
\begin{document}
\maketitle
\noindent\textit{Editorial note.} All mathematical statement and proof text is reproduced verbatim from the author's chapter etale.tex of the Stacks Project, fixed at commit a04446e57ec1fbc252a871afcec7752fb2807b14. The pinned preamble is supplied as its own file; the wrapper class, the xr-hyper wiring and the \textbackslash{}externaldocument list are omitted because the assembler supplies the class.\par\medskip
\noindent\textit{Editorial note.} Numbering inside this standalone excerpt is editorial, not the full Stacks book print numbering.\par\medskip
\noindent\textit{Editorial note.} The candidate is from the corrected band-I follow-on screen (line C) of the extractor's stacks-candidates line: its author label is absent from the live master in every form searched and it belongs to none of the earlier judged pools. The H2 judgement and its reason were made in this module by reading the author's own statement and complete proof in the fixed checkout. Original candidate key: stacks:etale.tex:lemma-fully-faithful.\par\medskip
\noindent\textit{Editorial note.} Reference labels in the proof resolve to navigation pointers to the original author source; these pointers do not reproduce proofs of the prerequisites.\par\medskip
\noindent\textit{Allowed standard prerequisites (editorial).}\begin{itemize}
\item Prerequisites named by the author inside the counted span. The counted span is reproduced verbatim, including the author's own citations. Its prerequisites are the earlier results of the same fixed Stacks Project version that the author cites there and in the chapter's lead-in.
\end{itemize}\medskip
\ifdefined\hypertarget\hypertarget{sourcepointer-1}{}\fi
\noindent\textit{Source reference.} Original equation 0BTI: sources/etale.tex, source lines 2035--2035. Navigation pointer used in the counted span; the span contains the pointer only and the cited result is not reproduced inside it. Registered as a non-load-bearing navigation stub in the established form for Stacks material. No bibliography citation is involved.\par\medskip
\ifdefined\hypertarget\hypertarget{sourcepointer-2}{}\fi
\noindent\textit{Source reference.} Original lemma 0BTJ: sources/etale.tex, source lines 2045--2045. Navigation pointer used in the counted span; the span contains the pointer only and the cited result is not reproduced inside it. Registered as a non-load-bearing navigation stub in the established form for Stacks material. No bibliography citation is involved.\par\medskip
\ifdefined\hypertarget\hypertarget{sourcepointer-3}{}\fi
\noindent\textit{Source reference.} Original lemma 01KV: sources/schemes.tex, source lines 4337--4337. Navigation pointer used in the counted span; the span contains the pointer only and the cited result is not reproduced inside it. Registered as a non-load-bearing navigation stub in the established form for Stacks material. No bibliography citation is involved.\par\medskip
\ifdefined\hypertarget\hypertarget{sourcepointer-4}{}\fi
\noindent\textit{Source reference.} Original lemma 01KS: sources/schemes.tex, source lines 4262--4262. Navigation pointer used in the counted span; the span contains the pointer only and the cited result is not reproduced inside it. Registered as a non-load-bearing navigation stub in the established form for Stacks material. No bibliography citation is involved.\par\medskip
\ifdefined\hypertarget\hypertarget{sourcepointer-5}{}\fi
\noindent\textit{Source reference.} Original proposition 024T: sources/etale.tex, source lines 467--467. Navigation pointer used in the counted span; the span contains the pointer only and the cited result is not reproduced inside it. Registered as a non-load-bearing navigation stub in the established form for Stacks material. No bibliography citation is involved.\par\medskip
\ifdefined\hypertarget\hypertarget{sourcepointer-6}{}\fi
\noindent\textit{Source reference.} Original theorem 025G: sources/etale.tex, source lines 1370--1370. Navigation pointer used in the counted span; the span contains the pointer only and the cited result is not reproduced inside it. Registered as a non-load-bearing navigation stub in the established form for Stacks material. No bibliography citation is involved.\par\medskip
\begin{lemma}
\label{lemma-fully-faithful}
Assume $f : X \to S$ is submersive and any \'etale base change
of $f$ is submersive. Then the functor
(\ref{equation-descent-etale}) is fully faithful.
\end{lemma}

\begin{proof}
By Lemma \ref{lemma-faithful} the functor is faithful.
Let $U_1 \to S$ and $U_2 \to S$ be \'etale morphisms
and let $a : X \times_S U_1 \to X \times_S U_2$ be a
morphism compatible with canonical descent data.
We will prove that $a$ is the base change of a morphism $U_1 \to U_2$.

\medskip\noindent
Let $U'_2 \subset U_2$ be an open subscheme. Consider
$W = a^{-1}(X \times_S U'_2)$. This is an open subscheme
of $X \times_S U_1$ which is compatible with the canonical
descent datum on $V_1 = X \times_S U_1$. This means that the
two inverse images of $W$ by the projections
$V_1 \times_{U_1} V_1 \to V_1$ agree. Since $V_1 \to U_1$
is surjective (as the base change of $X \to S$) we conclude
that $W$ is the inverse image of some subset $U'_1 \subset U_1$.
Since $W$ is open, our assumption on $f$ implies that $U'_1 \subset U_1$
is open.

\medskip\noindent
Let $U_2 = \bigcup U_{2, i}$ be an affine open covering.
By the result of the preceding paragraph we obtain an open
covering $U_1 = \bigcup U_{1, i}$ such that
$X \times_S U_{1, i} = a^{-1}(X \times_S U_{2, i})$.
If we can prove there exists a morphism $U_{1, i} \to U_{2, i}$
whose base change is the morphism
$a_i : X \times_S U_{1, i} \to X \times_S U_{2, i}$
then we can glue these morphisms to a morphism $U_1 \to U_2$
(using faithfulness). In this way we reduce to the case that
$U_2$ is affine. In particular $U_2 \to S$ is separated
(Schemes, Lemma \ref{schemes-lemma-compose-after-separated}).

\medskip\noindent
Assume $U_2 \to S$ is separated. Then the graph $\Gamma_a$ of $a$
is a closed subscheme of
$$
V = (X \times_S U_1) \times_X (X \times_S U_2) = X \times_S U_1 \times_S U_2
$$
by Schemes, Lemma \ref{schemes-lemma-semi-diagonal}.
On the other hand the graph is open for example
because it is a section of an \'etale morphism
(Proposition \ref{proposition-properties-sections}).
Since $a$ is a morphism of descent data, the two inverse images of
$\Gamma_a \subset V$ under the projections
$V \times_{U_1 \times_S U_2} V \to V$ are the same.
Hence arguing as in the second paragraph of the proof we
find an open and closed subscheme $\Gamma \subset U_1 \times_S U_2$
whose base change to $X$ gives $\Gamma_a$. Then
$\Gamma \to U_1$ is an \'etale morphism whose base change
to $X$ is an isomorphism. This means that $\Gamma \to U_1$
is universally bijective, hence an isomorphism
by Theorem \ref{theorem-etale-radicial-open}.
Thus $\Gamma$ is the graph of a morphism $U_1 \to U_2$
and the base change of this morphism is $a$ as desired.
\end{proof}

\end{document}
